\chapter[Partial Differential Equations]{Partial Differential Equations}

\section[Lecture 1 (02-03) -- {What are the PDEs?}]{Lecture 1 (02-03)Initial and Boundary Conditions}

\subsection{What are the PDEs?}
\begin{align}
    \label{eq1} \pdv{u}{t}=\pdv[2]{u}{x}+\pdv[2]{u}{y}-u 
    \\ u=u(x,y,t)
\end{align}
\eqref{eq1} is a second order PDE
 

\subsection{Classical solutions}
To simplify, we consider here equations of one unknown function $u$ of one or more variable
\\By a solution we mean a sufficiently smooth function $u$ that satisfies the PDE at every point of the domain of its definition($D$).
\\Sometimes, we will have to impose some boundary conditions on the domain $D$.
\\$C^n=\{u:D \rightarrow R:$ n-times continuous differentiable w.r.t all variables  \} $n$ is the order of the equation
\begin{example}{1}{}
    \begin{itemize}
        \item Heat equation: \[\pdv{u}{t}=\pdv[2]{u}{x}\]
        \[u=u(x,t)\]
        \[(t,x) \in D \subset R^2\]
        \[u(t,x), \pdv{u}{t}(t,x),\pdv{u}{x}(t,x),\pdv[2]{u}{x}(t,x),\pdv[2]{u}{t}(t,x)\]
        \[n=2\]
        \[\pdv{u}{x}{t}=\pdv{u}{t}{x} \in C^2\]
        \[u \in C^2(D)\]
        \item Some Solution:
        \[u(t,x)=t+\frac{x^2}{2},D=R^2\]
        \[u(t,x)=\frac{e^{-\frac{x^2}{4t}}}{2\sqrt{\pi t}},D=\{(t,x):t>0\}\]
        \[u(t,x)=e^{-t+ix}=e^{-t}cos(x)+ie^{-t}sin(x),D=R^2\]
    \end{itemize}
\end{example}
\subsection{Initial and Boundary Conditions}
What is Initial data?
\\$u(0,x)=u_0(x)$
\\
\\What are Boundary Conditions?
rigurous definition of Boundary
\begin{itemize}
    \item Dirichlet Boundary Condition: $u(t,x)|_{\partial{D}}=g(t,x)$
    \item Neumann Boundary Condition: When the normal derivative of $u$ is specified $\pdv{u}{n}$ where $n$ is the normal to the boundary $\partial(D)$
    \[\frac{\partial u}{\partial n}(t,x) = g(t,x)\]
    \item Mixed Boundary Condition: Combination of Dirichlet and Neumann
\end{itemize}
\subsection{Linear and Nonlinear waves}
\subsubsection{Stationary waves}
\begin{align}
    \label{eq2} \pdv{u}{t}&=0
    \\u&=u(t,x)
\end{align}
This is a first order linear homogeneuous PDE
\\Solution of the form $u=u(t)=c, \forall c \in R$
\\Let's integral \eqref{eq2} :
\[ 0=\int_{0}^{t}\frac{\partial u}{\partial t}(s,x)ds=u(t,x)-u(0,x)\]
\\Thus if an initial conditional
\[u(0,x)=f(x)\]
\[u(t,x)=f(x), \forall t,x\] 
Thus for any given $f \in C^1(R)$, we have a solution of \eqref{eq2} and this is a stationary wave
\\The domain of this solution is $R$ if $f$ is defined everywhere on $R$
\\Power of domain to change the solution:
\\Consider the ODE:
$$\dv{u}{t}=0, D=(-\infty,0)\cup(0,\infty) $$
\\The solution is:
\[
    u(t)=\begin{cases}
        c_1, t<0
        \\c_2, t>0
    \end{cases}
    \forall c_1,c_2 \in R
\]
\\similarly, for the PDE in example \eqref{eq2}:
$$
    u(t,x)=\begin{cases}
        0, x>0
        \\x^2, x\leq0,t>0
        \\-x^2, x\leq0,t<0
    \end{cases}
$$ 
it is a $ \mathbb{C}^1 $ solution of \eqref{eq2} on $ \mathbb{R}^2 $ and it is a classical solution.

\section[Lecture 2 (02-05) -- {Transport and Traveling waves}]{Lecture 1 (02-05)Transport and Traveling waves}
\subsection{uniform transport}
\begin{equation}
    \label{eq4} \pdv{u}{t}+c\pdv{u}{x}=0
\end{equation}
\begin{equation}
    \label{eq5} u=u(0,x)=f(x)
\end{equation}
$$
    x\in \mathbb{R},t_0=0,t-t_0=t,c\in \mathbb{R}
$$ 
\\\textbf{use change of variable:}  
$$
    \xi =x-ct
$$ 
$ \xi $ is called characteristic variable
\\let's look for a solution of \eqref{eq4} under the form $$ u(t,x)=v(t,\xi)=v(t,x-ct) $$ 
\textbf{use chain rule:}
$$
    \pdv{u}{t}=\pdv{v}{t}-c\pdv{v}{\xi}
$$ 
$$
    \pdv{u}{x}=\pdv{v}{\xi}
$$ 
Thus:
$$
    0=\pdv{u}{t}+c\pdv{u}{x}=\pdv{v}{t}-c\pdv{v}{\xi}+c\pdv{v}{\xi}= \pdv{v}{t}
$$ 
\textbf{Conclusion:}$$
    \pdv{v}{t}=0
$$ 
Solution is constant $ v(\xi) $
$$
u(t,x)=v(\xi)=v(x-ct)
$$  
\begin{proposition}{}{}
if u solves \eqref{eq4}, then u has the form $u(t,x)=v(x-ct)$ where v is a function of one variable that is $ \mathbb{C}^1 $-smooth.
\\Acutually, v is satisfying $$
    u(0,x)=v(x)=f(x)
$$ 
\end{proposition}
\textbf{Conclusion:}
$$\begin{cases}
    \pdv{u}{t}+c\pdv{u}{x}=0
    \\ 
    u=u(0,x)=f(x)
\end{cases}$$
The solution is $$ u(t,x)=f(x-ct) $$ 
\\$ f(x-ct) $ is a traveling wave solution moving at velocity $ c $:
\\if c>0, the wave moves to the right
\\if c<0, the wave moves to the left
\\if c=0, the wave is stationary 
\begin{example}{}{}
$$\begin{cases}
    \pdv{u}{t}+2\pdv{u}{x}=0
    \\ 
    u=u(0,x)=\frac{1}{1+x^2}
\end{cases}$$
The solution is $$ u(t,x)=\frac{1}{1+(x-2t)^2} $$
The lines $ x=ct+k $ are called the characteristic lines for the equation \eqref{eq4}
\\on the characteristic lines, the solution is always constant
\\on geometric interpretation, $u(t,x)=f(x-ct)=f(k)$
\end{example}
Consider the equation:
\begin{equation}
    \label{eq6} \pdv{u}{t}+c\pdv{u}{x}+au=0
\end{equation}
where a and c are given numbers.
\\\textbf{Use change of variable:}
$$
    u(t,x)=v(t,\xi)=v(t,x-ct)
$$ we will get by the same computation
$$
    \pdv{v}{t}+av=0
$$ 
\\The solution will be$$
    v(\xi)=e^{-at}f(\xi)
$$ 
$$
    u(t,x)=v(t,\xi)=v(x-ct)=e^{-at}f(x-ct)
$$ 
\\So,
$$
    \begin{cases}u(t,x)=e^{-at}f(x-ct)\\u(0,x)=f(x)\end{cases}
$$ 
\subsection{Nonuniform transport}
\begin{equation}
    \label{eq7} \pdv{u}{t}+c(x)\pdv{u}{x}=0
\end{equation}
\\In the uniform case(c constant), characteristic line is x=ct+k  (the solution of ODE below)
$$
    \dv{x}{t}=c
$$ 
\\In the general non-uniform case, we can still define characteristic curves
$$
    \frac{d}{dt}x(t)=c(x(t))
$$ 
\\Let's show that the solutions of \eqref{eq7} are constants along the characteristic curves (t,x(t))
\\We need to check that $ u(t,x(t))=h(t) $ is constant of t 
\\compute the derivative of $ h(t)$ 
\begin{align*}
    \frac{d}{dt}h(t) &= \frac{d}{dt}u(t,x(t)) = \pdv{u}{t}(t,x(t))+\pdv{u}{x}(t,x(t))\cdot\dv{x}{t} \\
    &= \pdv{u}{t}(t,x(t))+c(x(t)) \cdot \pdv{u}{x}(t,x(t)) = 0
\end{align*}
h(t) and thus u(t,x(t)) is constant along the characteristic curves
\begin{definition}{}{}
The graph of a solutioni of the ODE $$\frac{d}{dt}x(t)=c(x(t))$$ is called a characteristic curve for the transport equation \eqref{eq7} with speed c(x)
\end{definition}
Thus at each point (t,x) the slope of the characteristic curve is c(x) (the wave speed)
\\In the particular case c(x)=c, the characteristic curves are straight lines of slope c
\\Moreover,
\\ \eqref{eq7} is an autonomous first order ODE which can be "solved" using the method of separation of variables
\\Assume that c(x)$ \neq $0 
\begin{align*}
    \frac{dx(t)}{dt}&=c(x(t))
\\\frac{dx}{c(X)}&=dt \rightarrow
\\ \beta (x)&=\int\frac{dx}{c(x)}=t+k
\end{align*}
\\Thus characteristic curves satisfies the implicit equatioin$$\beta(x)=t+k$$
\\We can find x from the equality using $ \beta^{-1} $(the inverse of $ \beta $) 
\begin{align*}
    \beta(x)&=t+k
    \\ \beta^{-1}\beta(y)&=y
    \\x(t)&=\beta^{-1}(t+k)
\end{align*}
\begin{remark}{}{}
How to determine where $c(x_0)=0$?
\\Solve this ODE$$
    \begin{cases}
        \frac{dx}{dt}=c(x)
        \\x(t_0)=x_0
    \end{cases}
$$ 
and $ x(t)=x_0 $ wil be the characteristic curve going through ($t_0,x_0$) 
\end{remark}
From the proposition, we know that $ u(t,x(t)) $=constant
\begin{align*}
    u(t,x(t))&=v(k)
    \\&=v(\beta(x(t))-t)
\end{align*} 
\\where k is the constant determining the characteristic curve (t,x(t))
\begin{align*}
    u(t,x)=v(\xi)
    \\ \text{where    } \xi=\beta(x)-t
\end{align*}
\\which we will call characteristic variable
\begin{equation}
    \label{eq8} u(t,x)=v(\beta(x)-t)
\end{equation}
\\This formula way fail at points where c vanishes
\begin{align*}
    t&=0
    \\u(0,x)=f(x)&=v(\beta(x)-0)
    \\f(x)&=v(\beta(x))
    \\v(\xi)&=f(\beta^{-1}(\xi))
    \\&=f\cdot \beta^{-1}(\xi)
    \\u(t,x)&=f\cdot\beta^{-1}(\beta(x)-t)
\end{align*}
\\This is the formula for the solution of  \eqref{eq7} in the general case
\\which is:$$\begin{cases}
\pdv{u}{t}+c(x)\pdv{u}{x}=0
\\u(0,x)=f(x)
\end{cases}
$$ 